%!TEX root = tesis.tex

La gestión de residuos agroindustriales constituye un problema ambiental estratégico
por sus implicaciones sobre el cambio climático, la calidad ambiental y la
sostenibilidad de los sistemas productivos. A escala global, una fracción
significativa de la biomasa residual generada en las cadenas agroalimentarias no
se reincorpora de manera eficiente a circuitos de aprovechamiento, sino que es
dispuesta o transformada bajo condiciones que favorecen pérdidas de carbono y
emisiones de gases de efecto invernadero (GEI). En este contexto, la fracción
orgánica de los residuos adquiere especial relevancia, dado que su degradación
biológica puede generar emisiones de dióxido de carbono (\ce{CO2}) y, bajo
condiciones anaerobias, de metano (\ce{CH4}), un gas con mayor potencial de
calentamiento global que el \ce{CO2}.

En Colombia, esta problemática resulta particularmente relevante por la relación
entre la producción agropecuaria, la generación de residuos orgánicos y la
estructura nacional de emisiones. El Inventario Nacional de Emisiones y
Absorciones Atmosféricas presentado por el IDEAM muestra que las categorías
asociadas a agricultura, silvicultura y otros usos de la tierra, junto con
aquellas vinculadas a residuos, mantienen un peso importante dentro del perfil
nacional de GEI, lo que refuerza la necesidad de mejorar la caracterización de
fuentes y procesos con alta variabilidad e incertidumbre
\parencite{ideam2021inventario}.
Esta necesidad es consistente, además, con la Contribución Determinada a Nivel
Nacional (NDC) actualizada de Colombia, en la cual el país ratificó la meta de
reducir en 51\% sus emisiones proyectadas al año 2030 y reconoció la economía
circular como una estrategia relevante para la mitigación
\parencite{tiriaACTUALIZACIONCONTRIBUCIONDETERMINADA}.

Desde esta perspectiva, el tratamiento de residuos orgánicos no puede evaluarse
únicamente en función de la reducción de volumen, la estabilización del material
o la obtención de subproductos. También debe valorarse a partir de su desempeño
ambiental real, particularmente de las emisiones generadas durante el proceso.
Métodos convencionales como el compostaje abierto o la disposición en rellenos
sanitarios presentan limitaciones importantes: el primero puede emitir
\ce{CH4} y \ce{N2O} cuando no se mantienen condiciones adecuadas de aireación,
humedad y temperatura, mientras que el segundo favorece procesos anaerobios
ampliamente reconocidos como fuente de metano. En consecuencia, la selección de
tecnologías de tratamiento exige incorporar criterios de eficiencia climática,
además de los criterios operativos tradicionales.

En este escenario, la bioconversión de residuos mediante \textit{Hermetia
illucens} ha emergido como una alternativa promisoria dentro de enfoques de
economía circular y valorización de subproductos orgánicos. Su interés radica en
la capacidad de transformar sustratos residuales en biomasa de valor agregado y
en materiales potencialmente aprovechables en tiempos inferiores a los de otras
rutas biológicas convencionales. No obstante, la sostenibilidad de esta
tecnología no debe asumirse de forma automática. Su desempeño ambiental depende
de las condiciones del sustrato, de la interacción entre larvas y microbiota, y
de la dinámica de oxigenación del medio, factores que pueden inducir transiciones
entre estados predominantemente aerobios y anaerobios con efectos directos sobre
las emisiones de GEI.

Uno de los principales vacíos en este campo es que la evaluación de emisiones en
procesos de bioconversión suele apoyarse en mediciones puntuales, campañas
discontinuas o factores de emisión genéricos. Este tipo de aproximaciones resulta
insuficiente para describir procesos biológicos altamente dinámicos, en los que
las emisiones cambian con el tiempo en función de la humedad, la temperatura, la
densidad larval, la composición del sustrato y el estado metabólico del sistema.
Como resultado, persisten incertidumbres que limitan la comparación rigurosa
entre tecnologías, la identificación de eventos críticos de emisión y la
incorporación de información robusta en inventarios, análisis ambientales y
estrategias de mitigación.

Frente a esta brecha, la integración de sensórica, modelado matemático y
aprendizaje automático ofrece una oportunidad metodológica para desarrollar
sistemas híbridos de monitoreo capaces de estimar emisiones de manera continua,
dinámica y operativamente útil. Este enfoque resulta especialmente pertinente en
la bioconversión con \textit{H. illucens}, donde las emisiones emergen de
interacciones no lineales entre el sustrato, la biomasa, las condiciones
microambientales y el tiempo de residencia. En este sentido, más que registrar
concentraciones instantáneas, el reto científico consiste en construir
herramientas que permitan interpretar la dinámica del proceso y fortalecer la
gestión ambiental de esta tecnología.

En consecuencia, esta investigación se ubica en la intersección entre gestión de
residuos agroindustriales, cambio climático, monitoreo ambiental y modelación
computacional, ver figura \ref{fig:posicionamiento_tesis}. Su pertinencia radica en que responde a una necesidad concreta del
contexto colombiano: generar herramientas de estimación dinámica que contribuyan
a evaluar con mayor precisión el desempeño ambiental de la bioconversión de
residuos con \textit{Hermetia illucens}, y aportar evidencia útil para la
mitigación de GEI y la transición hacia sistemas agroindustriales más
sostenibles.

\shorthandoff{>}
\shorthandoff{<}

\begin{figure}[htbp]
\centering
\begin{tikzpicture}[
    every node/.style={font=\small},
    area/.style={
        ellipse,
        draw,
        thick,
        minimum width=4.2cm,
        minimum height=2.2cm,
        fill opacity=0.18,
        text opacity=1,
        align=center,
        font=\small\bfseries
    },
    thesis/.style={
        rectangle,
        rounded corners=6pt,
        draw=black!80,
        very thick,
        fill=white,
        text width=4.0cm,
        align=center,
        font=\footnotesize,
        inner sep=6pt
    }
]

% Cuatro elipses en las esquinas
\node[area, fill=blue!40] (A) at (-3.0,  2.4)
    {Gestión de\\residuos\\agroindustriales};

\node[area, fill=green!40] (B) at ( 3.0,  2.4)
    {Cambio\\climático y\\emisiones GEI};

\node[area, fill=orange!40] (C) at (-3.0, -2.4)
    {Monitoreo\\ambiental y\\sensórica};

\node[area, fill=purple!40] (D) at ( 3.0, -2.4)
    {Modelación\\computacional\\e IA};

% Flechas desde las áreas hacia el centro
\draw[->, thick, blue!60]   (A) -- (0, 0);
\draw[->, thick, green!60]  (B) -- (0, 0);
\draw[->, thick, orange!60] (C) -- (0, 0);
\draw[->, thick, purple!60] (D) -- (0, 0);

% Nodo central: la tesis
\node[thesis] (T) at (0, 0) {
    \textbf{Tesis}\\[4pt]
    Estimación dinámica\\de emisiones en\\
    bioconversión con\\
    \textit{Hermetia illucens}\\[4pt]
    \textcolor{gray}{(contexto colombiano)}
};

\end{tikzpicture}
\caption{Posicionamiento de la investigación en la intersección de cuatro
áreas temáticas: gestión de residuos agroindustriales, cambio climático,
monitoreo ambiental y modelación computacional.}
\label{fig:posicionamiento_tesis}
\end{figure}

\shorthandon{<}
\shorthandon{>}
